% chapters/abbreviations.tex
\addcontentsline{toc}{section}{ПЕРЕЛІК УМОВНИХ СКОРОЧЕНЬ, СИМВОЛІВ І ТЕРМІНІВ}
{\centering\bfseries ПЕРЕЛІК УМОВНИХ СКОРОЧЕНЬ, СИМВОЛІВ І ТЕРМІНІВ\par}
\vspace{1em}

\begin{tabular}{p{2.5cm}p{11.5cm}}
  АЦП    & аналого-цифровий перетворювач \\
  БКР    & бакалаврська кваліфікаційна робота \\
  ДПФ    & дискретне перетворення Фур'є \\
  ЕЕГ    & електроенцефалограма \\
  КС     & керуюче слово \\
  МРТ    & магнітно-резонансна томографія \\
  ОП     & оперативна пам'ять \\
  ПЗ     & програмне забезпечення \\
  ПК     & персональний комп'ютер \\
  ЦОС    & цифрове оброблення сигналів \\
  ШПФ    & швидке перетворення Фур'є \\
  \hline
  API    & Application Programming Interface (прикладний програмний інтерфейс) \\
  C4     & Context--Container--Component--Code (нотація опису архітектури) \\
  CI/CD  & Continuous Integration / Continuous Delivery \\
  DAG    & Directed Acyclic Graph (спрямований ациклічний граф) \\
  DFG    & Data Flow Graph (граф потоку даних) \\
  DFT    & Discrete Fourier Transform \\
  DSP    & Digital Signal Processing \\
  EMA    & Exponential Moving Average (експоненційне ковзне середнє) \\
  F1--F9 & функціональні вимоги до системи \\
  FFT    & Fast Fourier Transform \\
  FIFO   & First In, First Out \\
  FPGA   & Field-Programmable Gate Array \\
  GDPR   & General Data Protection Regulation \\
  IEEE   & Institute of Electrical and Electronics Engineers \\
  IEEE 754 & стандарт подання чисел з плаваючою крапкою \\
  LMS    & Learning Management System (система управління навчанням) \\
  LTI    & Learning Tools Interoperability \\
  MIT    & Massachusetts Institute of Technology (ліцензія з відкритим кодом) \\
  NF1--NF8 & нефункціональні вимоги до системи \\
  OFDM   & Orthogonal Frequency-Division Multiplexing \\
  RAM    & Random Access Memory \\
  rAF    & requestAnimationFrame \\
  RTL    & Register-Transfer Level (рівень регістрових передач) \\
  SDF    & Synchronous Data Flow (синхронний потік даних) \\
  SPA    & Single Page Application (односторінковий застосунок) \\
  SSR    & Server-Side Rendering \\
  TF     & Twiddle Factor (поворотний множник) \\
  UC1--UC9 & Use Cases (сценарії використання) \\
  UI     & User Interface (інтерфейс користувача) \\
  UML    & Unified Modeling Language \\
  UX     & User Experience (користувацький досвід) \\
  XSS    & Cross-Site Scripting \\
\end{tabular}
